\section{A directly verified selector--acceleration composition}
\label{sec:selector-fast-composite}

Two further 194-state tables are obtained without altering their parents'
source command protocols. The parents are the exact incoming AX11H and
LC\_AX11H tables, respectively. Write $M$ for either parent. First apply a
single history selector to controls $127,124$, producing $S$; then apply the
two incoming three-step accelerations to $S$, producing $N$. The result is
called PA194AX11H\_SF or PA194LC\_AX11H\_SF. The components are attributed
prior work; the new contribution here is their checked interaction on these
specific parents and the resulting graph-distinct literal constructions.

\paragraph{Why a checked interaction is necessary.}
The primary PA194 selector packet changes two pairs, $(12,38)$ and
$(127,124)$. The first pair cannot be copied into AX11H: its common read-one
row premise fails. In AX11H, $12/1=1R20$, whereas $38/1=1R122$.
The present construction uses only the second pair and verifies every
structural hypothesis directly from both literal parent tables.

\subsection{Complete changed rows}

All labels in the following table are native labels, with start 74 and
halt 254. The dense tables have start 0, working labels $0,\ldots,193$, and
halt 194. Each parent-to-child comparison changes exactly six rows:
\[
\begin{array}{c|c|c}
\text{row}&\text{parent}&\text{new}\\ \hline
104/1&1R124&1R127\\
127/0&0R14&0L124\\
124/0&0R62&0R8\\
124/1&1R138&1R59\\
26/1&1R42&1R91\\
101/0&0L26&0R17
\end{array}
\]
There are no other changed rows. Every one of these six rows is used during
the new blank initialization. Both resulting transition graphs have all
194 working controls reachable from their initial control.

\subsection{Selector correspondence}

Define $f(124)=127$ and let $f$ fix all other labels. In the parent,
\[
 \delta(127,0)=0R14,\quad \delta(124,0)=0R62,\qquad
 \delta(127,1)=\delta(124,1)=1R138.
\]
The helper identities are $\delta(8,0)=0R14$ and
$\delta(59,0)=0R62$. The complete table inspection shows that every
transition into 127 writes zero and moves right, and every transition into
124 writes one and moves right. The respective predecessor rows are
$110/0$ and $104/1$. Thus a run beginning in control 74 satisfies
\[
 (q=127\Rightarrow T(h-1)=0)\ \land\
 (q=124\Rightarrow T(h-1)=1).
\]
This is an entry-history invariant valid on every initial tape.

At ordinary checkpoints map $(q,h,T)$ to $(f(q),h,T)$. Outside the two
paired controls the child reproduces one parent step, with the target
mapped by $f$. The common read-one row likewise takes one step. On read
zero, the merged control moves left into selector 124; the selector reads
and preserves the old history bit, returns right to the original zero
through helper 8 or 59, and the helper performs the appropriate parent
transition. The block takes exactly three child steps. Only cells $h-1,h$
are inspected; their final contents agree with the parent, as does the
whole uninspected exterior. The four local cases and every ordinary row
are checked separately in the independent verifier.

Let $c_t=\operatorname{run}_M(t)$ with absorbing halt semantics and let
$F(q,h,T)=(f(q),h,T)$. Define
\[
 \tau(0)=0,\qquad \tau(t+1)=\tau(t)+
 \begin{cases}3,&c_t\text{ is in }124\text{ or }127\text{ and reads }0,\\
 1,&\text{otherwise}.\end{cases}
\]
Quantifier-free induction gives
\[
 \operatorname{run}_S(\tau(t))=F(c_t),\qquad
 t\leq\tau(t)\leq3t.
\]
Because $f$ maps no working state to halt, absorbing halting by a bounded
clock obeys
\[
 H_M(t)\Longrightarrow H_S(3t),\qquad
 H_S(t)\Longrightarrow H_M(t).                 \tag{SF1}
\]
The second implication follows by evaluating the displayed clock identity
at $t$: $S$ remains halted at time $\tau(t)\geq t$.

\subsection{Acceleration after the selector}

For the intermediate table $S$, the complete finite predecessor checks
still give
\[
 \operatorname{pred}(101)=\{96/0\mapsto0R101\},\qquad
 \operatorname{pred}(26)=\{11/0\mapsto0L26,\;101/0\mapsto0L26\},
\]
and every transition into 11 writes zero. Hence every $S$-run from 74 has
\[
\begin{split}
 q=101&\Rightarrow T(h-1)=0,\\
 q=11&\Rightarrow (T(h-1)=0\lor T(h+1)=0),\\
 q=26\land T(h)=1&\Rightarrow T(h+1)=T(h+2)=0.
\end{split}
\]
The proof is induction on one transition. For the last implication the
incoming transition writes the first zero; entry from 101 forces the new
read bit to zero, and entry from 11, when the new read bit is one, forces
the second zero. The selector does not invalidate these complete incoming
row conditions.

The two shortcuts are exact whole-tape identities:
\[
 26,1\;:\quad 1R42,\ 0R13,\ 0L91
       \quad\rightsquigarrow\quad 1R91,
\]
\[
 101,0\;:\quad 0L26,\ 0R194,\ 0R17
       \quad\rightsquigarrow\quad 0R17.
\]
Here native 194 is a working control, not the native halt label 254.
All intermediate controls are nonterminal. The independent verifier checks
all twelve admissible five-cell windows and verifies that no uninspected
cell is touched.

If $d_t=\operatorname{run}_N(t)$, define $\sigma(0)=0$ and increment
$\sigma$ by three precisely when $d_t$ executes one of the two shortcut
rows, otherwise by one. At the induction step, the induction hypothesis identifies the new boundary
with an actually reachable $S$-configuration. Hence the guards just proved
for $S$ apply to that boundary; they are not assumed separately for $N$.
Quantifier-free induction proves
\[
 \operatorname{run}_N(t)=\operatorname{run}_S(\sigma(t)),\qquad
 t\leq\sigma(t)\leq3t.
\]
Consequently
\[
 H_S(t)\Longrightarrow H_N(t),\qquad
 H_N(t)\Longrightarrow H_S(3t).                \tag{SF2}
\]
Combining (SF1) and (SF2) yields the particularly simple bounds
\[
 \boxed{H_M(t)\Longrightarrow H_N(3t),\qquad
        H_N(t)\Longrightarrow H_M(3t).}
\]
The factor is three in each direction, not nine: the selector and
acceleration slow the two simulations in opposite directions.

\subsection{Arithmetic scope}

Relative-to-head binary tape codes $(L,R)$ make every configuration
operation primitive recursive. A right move writing $w$ sends
$(L,2r+b)$ to $(2L+w,r)$; a left move, with $L=2\ell+a$, sends it to
$(\ell,4r+2w+a)$. The displayed entry invariants, run equalities, clock
bounds and halting persistence are quantifier-free predicates on explicitly
primitive-recursive functions. They therefore use only quantifier-free
induction. Pure intuitionistic logic now proves
\[
 \mathsf{iPRA}\vdash H_M\leftrightarrow H_N,
 \qquad
 \mathsf{iPRA}\vdash\neg H_M\leftrightarrow\neg H_N.
\]
This is a direct mathematical relative proof; it does not assume any source
calculus theorem or the consistency of PA. An absolute Con(PA) result is
obtained by composing it with the separately identified calibration of the
exact parent. No new end-to-end Lean kernel theorem is claimed.

\subsection{Exact initialization and novelty}

Both parent initializations take 57,922 steps to native control 84, head
$-2$, and the same complete seed frame. The selector branch is encountered
24 times, so the intermediate startup has $57,922+2\cdot24=57,970$ steps.
The two accelerations are used 20 and 6 times, giving
\[
 57,970-2(20+6)=57,918
\]
steps for each final machine. The independent program checks complete tape
support equality after every simulation block throughout these prefixes.
This concrete binding check is separate from the all-time induction proof.

The graph comparator compares full canonical tuples against all 139
spanning 194-state classes in the current pinned corpus and the additional
nonspanning diagnostic, allowing every spanning root, control renaming,
reflection, global symbol exchange and ignored halt-edge actions. Global
symbol exchange is an extra structural screen, not a zero-blank semantic
transformation. The two new graphs are mutually distinct and match none
of those prior tables. This establishes corpus-relative graph novelty;
it is not a claim of worldwide priority or of a 194-state lower bound.
